\documentclass[brazilian,11pt,a4paper]{article}
\usepackage[T1]{fontenc}
\usepackage[utf8]{inputenc}
\usepackage{lmodern}
\usepackage[brazilian]{babel}
\usepackage[margin=2.5cm]{geometry}
\usepackage{amsmath,amssymb}
\usepackage{booktabs,array,graphicx}
\usepackage{microtype}
\usepackage{enumitem}
\usepackage{xurl}
\usepackage[dvipsnames]{xcolor}
\usepackage[colorlinks=true,linkcolor=Maroon,urlcolor=Blue,citecolor=Blue]{hyperref}
\setlist{itemsep=2pt,topsep=4pt}
\setlength{\parskip}{4pt}
\setlength{\parindent}{0pt}

\hypersetup{
  pdftitle={Olimpo: classificação de notícias em português e plataforma educativa},
  pdfauthor={Equipe Olimpo, PUC-Campinas e Instituto de Pesquisas ELDORADO},
  pdflang={pt-BR}
}

\title{\textbf{Olimpo: o que um classificador linear aprende sobre notícias falsas em português, e como isso alimenta uma plataforma educativa}}
\author{Equipe Olimpo\\PUC-Campinas \,\&\, Instituto de Pesquisas ELDORADO}
\date{Outubro de 2026}

\begin{document}
\maketitle

\begin{abstract}
\noindent O Olimpo é uma plataforma gamificada que treina o reconhecimento de notícias falsas por meio de votação em salas e feedback socrático. Este artigo resume as conclusões da trilha de aprendizado de máquina que sustenta o produto e o que, ao final, está de fato em uso. No Fake.br-Corpus (7.200 notícias), um \textit{LinearSVC} sobre TF-IDF de palavras e caracteres atinge F1 $\approx$ 0,93 em validação cruzada, e nem redução de dimensionalidade, nem kernels, nem árvores, nem metadados morfossintáticos trazem ganho relevante. Em títulos do FakeRecogna, porém, o F1 cai para $\approx$ 0,61 e o AUC para 0,66--0,71: parte do sinal é artefato da fonte (autoria, tipografia, tamanho). Os métodos não supervisionados mostraram-se úteis para \emph{descrever} estilos de escrita, não para classificar. O modelo em produção é um \textit{LinearSVC} calibrado com atributos spaCy, $\chi^2$ e SVD, cuja pontuação se decompõe exatamente em n-gramas e atributos. Ele já roda em um serviço Python próprio, consumido pelo backend do jogo, e seu resultado é revelado só após o fechamento coletivo de cada rodada e tratado como \emph{indício}, não como veredito. Um segundo componente, baseado em regras do FP-Growth, mostra durante a leitura como a escrita do trecho se compara à de cada classe do corpus, sem classificar.

\medskip
\noindent\textbf{Palavras-chave:} desinformação, classificação de texto, SVM, FP-Growth, interpretabilidade, educação midiática.
\end{abstract}

\section{Introdução}

Plataformas digitais priorizam engajamento sobre precisão, e ferramentas de checagem que apenas carimbam um texto como ``falso'' tendem a ser rejeitadas por quem já discorda delas. A proposta do Olimpo é outra: em salas multiplayer ou no modo solo, os jogadores votam entre \emph{Fake}, \emph{Fato real} e \emph{Tenho dúvida}, e ao fim da rodada recebem o gabarito com uma justificativa pedagógica e perguntas no estilo socrático. O componente de IA fornece a classificação e os motivos usados nesse feedback.

A pergunta que guia o trabalho técnico é: \emph{o que um modelo treinado em um corpus rotulado de notícias consegue, de fato, afirmar sobre textos novos?} A resposta condiciona como o produto deve apresentar a saída do modelo.

\section{Dados e protocolo}

\textbf{Corpus.} Usamos o Fake.br-Corpus (NILC/USP): 7.200 notícias em português, 3.600 falsas e 3.600 verdadeiras. As falsas são bem mais curtas (mediana de 157 palavras contra 918). Para checar generalização, usamos o FakeRecogna (11.903 registros), restrito a \emph{títulos}: 2.484 falsos do domínio \texttt{boatos.org} e igual número de reais, com mediana de 13 palavras.

\textbf{Pré-processamento.} Normalização tipográfica (aspas, travessões, reticências, dígitos $\to$ 0), truncamento em 100 palavras e metadados em taxa por 100 palavras. Essas decisões vieram da análise exploratória: sem elas o modelo aprendia a fonte em vez do conteúdo. Partição estratificada 80/20 (5.760 treino, 1.440 teste) e validação cruzada de 5 folds com a métrica F1.

\textbf{Modelos.} Supervisionados: \textit{LinearSVC}, Regressão Logística (baseline interpretável) e Random Forest, sobre TF-IDF de palavras e caracteres, com e sem atributos morfossintáticos extraídos por spaCy. Não supervisionados: Isolation Forest, LOF, One-Class SVM, K-means, DBSCAN, HDBSCAN e mineração de regras de associação com FP-Growth.

\section{Resultados}

\subsection{Artefatos de fonte no corpus}

Três achados independentes mostram que parte do sinal vem da fonte, e não da linguagem da desinformação:
\begin{itemize}
  \item \textbf{Autoria:} 3.528 de 3.601 notícias sem autor são falsas e 3.527 de 3.599 com autor são reais (Cramér V = 0,960). O campo foi excluído dos modelos.
  \item \textbf{Formatação:} no texto bruto, espaços e quebras são 2,93\% dos tokens das falsas contra 0,38\% das verdadeiras. Após normalização e truncamento, caem a 0,00\% nas duas classes.
  \item \textbf{Convenções de redação:} marcadores temporais (``nesta'', ``quarta'') concentram-se nas verdadeiras e afetam a classificação.
\end{itemize}
Resta um sinal de tamanho: 17\% das falsas têm menos de 100 palavras, contra $\approx$0\% das verdadeiras.

\subsection{Aprendizado supervisionado}

\begin{table}[h]
\centering
\small
\begin{tabular}{@{}p{6.2cm}p{3.2cm}p{4.6cm}@{}}
\toprule
\textbf{Configuração} & \textbf{F1 (CV, 5 folds)} & \textbf{Observação}\\
\midrule
\textit{LinearSVC} word+char, $C{=}1$ & 0,9295 $\pm$ 0,0099 & referência interna\\
Regressão Logística word+char & 0,9182 ($C{=}1$) $\to$ 0,9281 ($C{=}10$) & baseline interpretável\\
Random Forest word+char & 0,9285 & protocolo anterior\\
Redução ($\chi^2$, L1, SVD) & --- & nenhuma supera a base\\
Kernels RBF/polinomial & --- & sem ganho\\
+ metadados spaCy ($w{=}0{,}03$) & 0,9344 & +0,005; ganho frágil\\
\textbf{spaCy + $\chi^2$ 10 mil + SVD(500)} & \textbf{0,9295} & gap 0,042 (base: 0,0705)\\
\bottomrule
\end{tabular}
\caption{Desempenho interno no Fake.br-Corpus.}
\end{table}

Em síntese: (i) em regime de muitos atributos e poucos documentos, o modelo linear simples é difícil de superar, e $C{=}1$ é o ponto ótimo; (ii) reduzir a dimensão só diminui a distância entre treino e validação (por ajuste menor), sem melhorar F1; (iii) os metadados do corpus tiveram efeito contraditório (--2,4 pt na Regressão Logística, +1,4 pt no SVM), e os atributos spaCy somam apenas +0,005 e só com peso baixo; (iv) no teste interno, spaCy e base empatam dentro do erro padrão ($\approx\pm$0,007).

\subsection{Generalização fora do corpus}

\begin{table}[h]
\centering
\small
\begin{tabular}{@{}llccc@{}}
\toprule
\textbf{Conjunto} & \textbf{Modelo} & \textbf{F1 macro} & \textbf{ROC-AUC} & \textbf{\% previsto fake}\\
\midrule
Interno (Fake.br) & base word+char & 0,9222 & 0,9830 & 49,7\\
Interno (Fake.br) & + spaCy & 0,9236 & 0,9846 & 49,6\\
Externo (FakeRecogna, títulos) & base word+char & 0,6119 & 0,7072 & 70,8\\
Externo (FakeRecogna, títulos) & + spaCy & 0,6134 & 0,6622 & 59,3\\
\bottomrule
\end{tabular}
\caption{Teste interno e teste externo (treino sempre no Fake.br).}
\end{table}

O desempenho não se transfere. O modelo tende a rotular como falso o que é curto e diferente do corpus, e as variáveis de maior peso (taxa e tamanho médio de sentenças) saem da distribuição de treino em títulos. Os atributos spaCy deslocam o limiar, mas \emph{pioram} a ordenação (AUC 0,707 $\to$ 0,662). Nenhuma técnica de dimensionalidade resolveu a distância de domínio e formato.

\subsection{Aprendizado não supervisionado}

Os detectores de novidade atingiram acurácias altas (Isolation Forest 0,97; DBSCAN 0,97; LOF 0,95), mas a leitura foi revisada: eles aprendem desvio em relação às notícias verdadeiras, e como a autoria separa quase por completo as classes, o resultado reflete o corpus, não a falsidade. Em holdout temporal a acurácia cai (Isolation Forest 0,86; DBSCAN 0,82). O agrupamento temático teve silhouette de 0,008 (K-means) e 0,040 (HDBSCAN), sem autorizar tratar grupos como classes de veracidade.

O modelo principal atual é o FP-Growth sem autoria da raiz de \path{machine-learning/unsupervised-learning/}: \texttt{fp\_growth\_principal.py} e \path{fp_growth_principal_sem_autoria.ipynb}, variante \texttt{sintaxe\_ampliada}. No run \path{fp-growth-metadados-ampliados-20261008T002311Z}, são 22 atributos ativos, 44 itens, 523 regras direcionais elegíveis e 345 padrões distintos. Na validação, 422 regras mantêm os filtros e 347 também têm redescoberta $\geq$80\% em 100 reamostragens dos grupos de treino. O teste conhecido é exploratório. O método descreve coocorrências e não classifica a veracidade de notícias.

Os quantis e regras usam apenas o treino canônico (4.320 notícias); validação e teste têm 1.440 notícias cada e preservam grupos de pares e duplicatas. A extração usa os primeiros 300 caracteres após normalização NFKC. O suporte mínimo é 0,08 e os filtros são confidence $\geq$0,50, lift $\geq$1,05 e Jaccard $\geq$0,10. Classe e autoria são usadas apenas na descrição posterior; o ranking de 20 candidatos não usa rótulos. Os relatórios estilísticos e o ranking anterior com 25 padrões são históricos e não representam este principal.

Na validação, 105 dos 345 padrões têm maior frequência dentro da classe Fake, 235 dentro de True e cinco empatam. Essas frequências não são probabilidades de falsidade. A união de todos os padrões cobre 1.439/1.440 notícias; a fila de revisão cobre 1.346, sem comprovar eficácia pedagógica. Ainda faltam avaliação externa e revisão editorial humana.

\section{O que está em uso}

A trilha de pesquisa chegou ao produto por uma camada de execução separada, o \texttt{model-engine}, que é consumida pelo servidor Next.js. Os dois componentes de IA têm funções, saídas e momentos de exibição distintos e propositalmente \emph{não} são combinados em um único escore.

\begin{table}[h]
\centering
\small
\begin{tabular}{@{}p{3.2cm}p{4.6cm}p{6.2cm}@{}}
\toprule
\textbf{Componente} & \textbf{Função} & \textbf{Onde aparece no jogo}\\
\midrule
Observação de escrita (FP-Growth) & Mostra padrões gramaticais presentes no trecho e a frequência da mesma combinação em cada classe do corpus & Painel ``Observe a escrita'', durante a leitura\\
Modelo supervisionado & Estima $P(\text{fake})$ com explicação por n-grama e atributo & ``Análise do modelo'', após todos votarem ou o tempo acabar\\
Gabarito cadastrado & Referência da pontuação dos jogadores & Veredito da rodada, identificado à parte da previsão\\
\bottomrule
\end{tabular}
\caption{Papéis dos componentes no jogo.}
\end{table}

\subsection{Modelo supervisionado}

\textbf{Modelo.} \textit{LinearSVC} ($C{=}1$, classes balanceadas) com 24 atributos spaCy (peso 0,03), seleção $\chi^2$ de 10 mil n-gramas e SVD de 500 dimensões, calibrado por sigmoide (versão \path{svm-spacy-chi2k10k-svd500-v1}). Iguala a base em F1 de validação cruzada (0,9295), reduz o gap de 0,0705 para 0,042 e tem F1 de 0,9243 no teste interno sem calibração. A versão calibrada registra F1 de 0,9208 e AUC de 0,9789 antes do reajuste final. A escolha se deve menos ao desempenho que à \emph{explicabilidade}: o escore se decompõe exatamente em contribuições por n-grama e por atributo de estilo, que viram os motivos exibidos. O artefato final foi reajustado em treino mais teste (7.200 notícias), de modo que as métricas internas citadas vêm de outro ajuste e não avaliam esse artefato de forma independente.

\textbf{Política de decisão.} O escore é $100\times P(\text{fake})$. A classificação é \emph{confiável} para $P\le0{,}35$, \emph{não confiável} para $P\ge0{,}65$ e \emph{incerta} entre os dois valores. Textos com menos de 30 palavras não são classificados, e a análise usa as primeiras 100 palavras do corpo extraído, como no treino. Os limiares são decisão de produto.

\textbf{Serviço.} O \texttt{model-engine} é um serviço HTTP em Python (\texttt{POST /supervised/analyze}), que recebe apenas o corpo da notícia e não registra seu conteúdo. Antes de carregar o modelo, confere o hash SHA-256 do artefato, o patch exato do Python (3.14.2) e as versões de scikit-learn, spaCy e demais bibliotecas; se algo divergir, somente o supervisionado fica indisponível, e as observações de escrita continuam funcionando. Numa medição local (20 análises aquecidas), a mediana foi de 140,6~ms e o p95 de 188,3~ms, com pico de 438~MiB de memória; com oito requisições simultâneas, três foram atendidas e cinco recusadas, pois a fila é intencionalmente curta (uma execução e duas vagas de espera). Essa medição não representa um teste de carga em produção.

\textbf{Integração no backend.} Um adaptador HTTP substituiu o \textit{mock} como implementação padrão do contrato de análise (\texttt{classification}, escore, \texttt{reasons}, \texttt{modelVersion}). A resposta é guardada em PostgreSQL com chave formada pelo conteúdo, pelo hash do artefato, pelo código de inferência e pela política, de modo que trocar a versão do modelo gera nova análise. Em um teste com quatro consultas concorrentes de duas instâncias do caso de uso houve uma única inferência e um único registro.

\textbf{Acesso e pontuação.} A rota \path{POST /api/news-prediction} exige sessão, participação na sala e rodada encerrada coletivamente: votar, avançar ou concluir individualmente não libera a previsão. Votos, placar e eventos em tempo real fecham sem chamar o modelo, e a pontuação continua vindo do gabarito cadastrado. Se o modelo falhar, a partida segue normalmente e apenas a previsão fica indisponível.

\textbf{Validação.} Um teste ponta a ponta com parser HTML real, modelo real, banco e Redis descartáveis e duas sessões de navegador verificou o bloqueio antes do fechamento, as cinco consultas concorrentes com um registro, a recarga, as telas desktop e celular, e o encerramento da partida com falha induzida do modelo. A suíte do web app reportou 419 testes aprovados, e a suíte Python 20, incluindo reprodução numérica das \textit{features} e das probabilidades do código original.

\subsection{Observação de escrita (FP-Growth)}

O FP-Growth foi levado ao produto como \emph{descrição}, não como classificador. O catálogo em uso vem da variante \texttt{sintaxe\_ampliada} (22 atributos, 523 regras direcionais elegíveis e 345 padrões distintos) e contém 20 padrões candidatos, selecionados sem usar classe ou autoria no ranking. Para um texto novo, o motor normaliza e recorta os primeiros 300 caracteres, extrai as medidas com limites congelados, sem recalcular quantis nem receber gabarito ou autor, e indica quais padrões correspondem \emph{integralmente}. Cada padrão traz uma explicação em linguagem simples (por exemplo, ``afirmações ligadas a um verbo, como \emph{disse que o prazo acabou}'') e a frequência da mesma combinação em cada classe da partição de validação (720 notícias falsas e 720 verdadeiras). Essas frequências são \emph{dentro de cada classe} do corpus, e não a probabilidade de a notícia ser falsa. O catálogo está marcado como \path{experimental_descriptive_comparison} e não passou por aprovação editorial humana.

\subsection{Plataforma}

Aplicação Next.js 16/React 19 em TypeScript, com PostgreSQL (Drizzle), Redis e eventos por SSE, organizada em Clean Architecture no backend e Feature-Sliced Design no frontend. Estão implementados autenticação, salas com PIN, votação em três níveis, ranking, desafios autorais, extração de notícias por URL com proteção contra SSRF e a interface por etapas. O ambiente completo (app, migrações, banco, Redis e \texttt{model-engine}) sobe com um único \texttt{docker compose}. A reflexão socrática exibida na interface ainda consiste em perguntas gerais fixas; a geração dinâmica de perguntas a partir dos motivos do modelo e dos padrões observados não existe.

\section{Discussão e limitações}

\begin{enumerate}
  \item \textbf{Dentro do corpus, o problema está resolvido por um modelo simples; fora dele, não.} F1 $\approx$ 0,93 contra $\approx$ 0,61 em títulos de outra fonte.
  \item \textbf{Rótulo e fonte estão confundidos.} Há um único corpus de treino, rotulado por origem, e o desempenho alto parcialmente mede reconhecimento da fonte.
  \item \textbf{Complexidade não comprou generalização.} Kernels, árvores, redução de dimensão e metadados não ajudaram; o ganho do modelo de produção é de explicabilidade e de menor ajuste ao treino.
  \item \textbf{Cuidados de avaliação.} As partições supervisionadas são aleatórias (pares temáticos podem cruzar treino e teste), e cada configuração externa foi avaliada uma única vez. As avaliações do FP-Growth descrevem o corpus, não medem generalização.
  \item \textbf{Sem avaliação em uso real.} O modelo funciona de ponta a ponta, mas nenhuma métrica foi medida em notícias submetidas por usuários, e a transferência externa é fraca (AUC $\approx$ 0,68 em títulos). A interface não deve exibir números sem origem medida, como a margem de erro e a dica de ``78\%'' que estavam fixas na tela.
  \item \textbf{Limites operacionais.} O marcador de fechamento da rodada no Redis expira em 24 horas; depois disso a previsão é bloqueada, mesmo em partidas finalizadas. A fila de inferência é curta e recusa picos de requisições simultâneas.
\end{enumerate}

Esses limites sustentam o posicionamento do produto: o Olimpo é uma plataforma \emph{educacional} que apresenta indícios, e não um detector definitivo de fake news; o modelo não verifica fatos.

\section{Conclusão}

O trabalho estabeleceu que um classificador linear sobre TF-IDF resolve o Fake.br-Corpus com F1 $\approx$ 0,93, que esse resultado é em boa parte específico do corpus e que técnicas mais sofisticadas não o melhoram nem o transferem. O valor mais robusto da trilha foi metodológico: identificar e neutralizar artefatos de fonte, separar descoberta de classificação e escolher um modelo cuja saída pode ser explicada ao jogador. O resultado prático é que a pesquisa chegou ao jogo: o modelo supervisionado roda em serviço próprio, com versão e hash verificados, e é revelado só depois do fechamento coletivo; o FP-Growth acompanha a leitura como descrição de escrita. Os próximos passos são: (1) medir o modelo com notícias reais extraídas pela plataforma; (2) coletar um conjunto externo de texto completo com fontes e épocas distintas e repetir a avaliação; (3) submeter o catálogo de observações a revisão editorial humana e a um estudo de compreensão com jogadores; (4) gerar dinamicamente as perguntas socráticas a partir dos motivos do modelo e dos padrões observados; e (5) persistir o fechamento da rodada de forma durável e dimensionar a fila de inferência para carga real.

\section*{Referências}
\begin{itemize}[leftmargin=*,itemsep=1pt]
  \item NILC/USP. \textit{Fake.br-Corpus}. Revisão fixa 780f5516c4ae070761632d98ac3368f3ded09d35.
  \item RECOGNA NLP. \textit{FakeRecogna}. \url{https://huggingface.co/datasets/recogna-nlp/FakeRecogna}.
  \item Lundberg, S.; Lee, S.-I. A unified approach to interpreting model predictions (SHAP). \url{https://shap.readthedocs.io/}.
  \item Scikit-learn. \textit{User guide}. \url{https://scikit-learn.org/stable/}.
  \item Equipe Olimpo. \textit{Documento técnico consolidado} (\path{Entregável Final - Completo}) e relatórios em \texttt{machine-learning/}, \texttt{model-engine/} e \path{docs/machine-learning/validacao-integracao-supervisionado.md} (branch \texttt{dev}), repositório \texttt{olimpo-fake-news-ai}, 2026.
\end{itemize}

\end{document}
